\section{Periodic-orbit catalogue and restricted operator construction}
\label{sec:orbits}

\subsection{Catalogue bridge}

The cycle ledger starts from the $a=6$ periodic catalogue through period twelve.
Each candidate was refined at high precision and classified relative to the
four state rectangles in \eqref{eq:hsets}.  The final classification contains
79 records inside the selected domain and 668 outside, with no unresolved or
root-failed record.  Theorem~\ref{thm:contraction} and
Lemma~\ref{lem:proof-period} make the admissible symbolic-necklace count exact
on $\Lam_*$.  In particular, M\"obius inversion gives
\begin{equation}
 p_n=\frac1n\sum_{d\mid n}\mu(d)\operatorname{tr}\!\left(A^{n/d}\right),
 \qquad
 (p_1,\ldots,p_{12})=(1,0,1,2,2,2,4,5,8,11,18,25),
 \label{eq:primitive-counts}
\end{equation}
whose entries sum to 79.  An
independent checker re-enumerates the primitive words, reclassifies the
persisted high-precision sequences, and verifies the trace/primitive-count
identity.  The numerical inside records agree period by period with the exact
symbolic necklaces.  The numerical catalogue therefore confirms the exact
count while supplying coordinates and multipliers not contained in the
symbolic enumeration.

For a periodic point $z_p$, the code stores the monodromy matrix $M_p$, its
determinant residual, the unstable multiplier modulus, and the time-reversal
partner when one exists.  The reciprocal relation is checked numerically as a
diagnostic of implementation quality; the exact determinant-one identity comes
from \eqref{eq:derivative-inverse}.

\subsection{Restricted weighted-cycle algebra and cutoff stability}

On the certified symbolic survivor, the full unweighted Euler determinant is
the exact finite-state identity
\begin{equation}
 D_{E,0}(z)=Z_0(z)^{-1}
 =\prod_p(1-z^{n_p})
 =\det(I-zA)=1-z-z^3-z^4,
 \label{eq:beta-zero-determinant}
\end{equation}
where the product is over primitive symbolic necklaces.  The stored degree-12
coefficient section reproduces this polynomial through degree twelve with zero
error at the stored precision.  For positive weights we
form the truncated Euler and periodic-point flat polynomials from the same
restricted orbit catalogue.  For either polynomial, the reported leading
modulus is $|z_*|^{-1}$, where $z_*$ is its smallest-modulus nonzero root.  The
period-cutoff sequence in Table~\ref{tab:r059-cycle-cutoff} shows stabilization of
the $\beta=1$ Euler and flat values near $0.29162845$.  The $\beta=1/2$
Euler value near $0.68284718$ is included only as a weight-sensitivity
diagnostic; it is not the target of the finite-volume comparison.

\begin{table}[H]
\centering
\caption{Restricted weighted-cycle cutoff sequence.  The orbit count is cumulative
through primitive period $N$.}
\label{tab:r059-cycle-cutoff}
\small
\begin{tabular}{rrrrr}
\toprule
$N$ & primitive orbits & $\lambda_{E,1/2}$ & $\lambda_{E,1}$ & $\lambda_F$\\
\midrule
8  & 17 & 0.682846853 & 0.291628304 & 0.291624276\\
9  & 25 & 0.682847077 & 0.291628328 & 0.291628474\\
10 & 36 & 0.682847204 & 0.291628476 & 0.291628512\\
11 & 54 & 0.682847184 & 0.291628450 & 0.291628441\\
12 & 79 & 0.682847185 & 0.291628451 & 0.291628452\\
\bottomrule
\end{tabular}
\end{table}

\begin{figure}[H]
\centering
\includegraphics[width=0.84\textwidth]{r059_cycle_convergence.pdf}
\caption{Restricted cycle-polynomial leading moduli versus period cutoff.  The
right panel resolves the agreement of the $\beta=1$ Euler and periodic-point
flat sections at the final cutoffs.  Because the flat determinant contains the
Euler determinant as its leading factor, this is a finite-section consistency
check rather than independent dynamical evidence.}
\label{fig:r059-cycle-convergence}
\end{figure}

At $N=12$, the two $\beta=1$ diagnostics differ by
$1.85\times10^{-10}$ in leading modulus.  In view of
\eqref{eq:flat-factorization}, their formal dependence motivates this empirical
finite-truncation and implementation check.  It neither implies equality of
the degree-12 zero sets nor bounds their root discrepancy.  The finite-volume
comparison below uses their common displayed $\beta=1$ value as a fixed cycle
reference; it does not identify the finite matrix with a continuous
thermodynamic operator.

\subsection{Absorbing finite-volume matrices}

Partition each state rectangle into an $m\times m$ array of cells $B_I$ and
write $\mu$ for planar Lebesgue measure.  The ideal open finite-volume matrix is
\begin{equation}
 (P_m)_{IJ}
 =\frac{\mu\!\left(B_I\cap H_6^{-1}(B_J)\right)}{\mu(B_I)},
 \qquad P_m\in\mathbb R^{4m^2\times4m^2}.
 \label{eq:ulam-definition}
\end{equation}
Rows index source cells and columns index target cells.  State-major ordering
is followed by local $y$ and local $x$ indices.  The implemented estimators
replace the normalized cell integral in \eqref{eq:ulam-definition} by either
tensor Gauss--Legendre weights, normalized to sum to one in each source cell,
or equal-weight scrambled Sobol samples.  A mapped point in $B_J$ contributes
to column $J$; a point outside $\mathcal N^\circ$, or exactly on a discarded
state-rectangle or local-cell boundary, contributes to escape.  The resulting
CSR matrices are nonnegative and substochastic up to the recorded
$10^{-12}$ audit tolerance.

The production panel uses
$m\in\{24,32,48,64,96,128\}$, Gauss--Legendre order eight, and two scrambled
Sobol estimators with 64 points per cell.  The independent schema checker
verifies matrix hashes, indexing, selected source rows, zero exact boundary
hits, and returned eigenpair residuals.  Table~\ref{tab:r059-cycle-operator}
shows the six finest-grid comparisons with the period-12 cycle reference
$\lambda_{\rm cyc}=0.29162845164474605$.

\begin{table}[H]
\centering
\caption{Finest-grid finite-volume comparison.  The relative gap is
$|\lambda_m-\lambda_{\rm cyc}|/|\lambda_{\rm cyc}|$.}
\label{tab:r059-cycle-operator}
\small
\begin{tabular}{lrrr}
\toprule
Estimator & $m$ & leading modulus $\lambda_m$ & relative gap\\
\midrule
Gauss $q=8$ & 96  & 0.291998537 & 0.1269\%\\
Gauss $q=8$ & 128 & 0.288096349 & 1.2112\%\\
Sobol 20260801 & 96  & 0.291078899 & 0.1884\%\\
Sobol 20260801 & 128 & 0.292758121 & 0.3874\%\\
Sobol 20260802 & 96  & 0.292764254 & 0.3895\%\\
Sobol 20260802 & 128 & 0.290990842 & 0.2186\%\\
\bottomrule
\end{tabular}
\end{table}

The median relative gap is $0.3030\%$.  The code applies ARPACK to the density
transpose $P_m^\top$, normalizes each returned Ritz vector, and records
\begin{equation}
 \eta_j=\norm{P_m^\top v_j-\lambda_jv_j}_2,
 \qquad \norm{v_j}_2=1.
 \label{eq:arpack-residual}
\end{equation}
Across these six records, all returned pairs have
$\eta_j\leq1.66\times10^{-11}$, verifying small algebraic residuals for those
Ritz pairs.  The first/second modulus separation is $22.8\%$--$29.2\%$ of the
first modulus, but only among the returned eigenvalues.  Here ``leading'' means
the largest modulus among the converged ARPACK eigenpairs requested with
\texttt{which=LM}.  For these nonnormal finite matrices, neither a small
residual nor separation within the returned set proves proximity to a true
eigenvalue or encloses the spectral radius.  The historical dyadic stability
gate remains negative because two Sobol trajectories are nonmonotone, so the
table records a common finite-resolution window rather than convergence.

\subsection{Variance and quadrature sensitivity}

An extended panel adds sixteen Sobol seeds and paired budgets 64 and 256 on all
six grids, together with Gauss orders 4, 8, and 12.  Increasing the budget reduces
cross-seed standard deviation at every grid; the six SD(256)/SD(64) ratios lie
between 0.250 and 0.512, with median 0.3207.  The strict paired-shift gate
still fails at $m=24$ and $m=32$, and the q=8/q=12 comparison is not stable on
all grids.  Thus sampling variance is material but does not explain the whole
historical failure.

The next section holds a finest-grid cloud fixed while changing the coarse
representation.  This supplies a paired sensitivity experiment rather than a
replacement of the preceding estimators.
